\begin{Solbook}{1}
 \par \noindent  a) $\dfrac {1}{9}$.\quad b) $0$.\quad c) $\dfrac {11}{18}$.\quad d) $9$.  \hfill \qedEX 
\end{Solbook}
\begin{Solbook}{2}
 \par \noindent  a) $x=-\dfrac {1}{5}$.\\ b) $x=\dfrac {4}{3}$.\\ c) $x=\dfrac {37}{20}$. Vì vế phải không âm nên $x\ge 0$; khi đó cả hai biểu thức trong dấu giá trị tuyệt đối đều dương.\\ d) $x=4$.  \hfill \qedEX 
\end{Solbook}
\begin{Solbook}{3}
 \par \noindent  a) $2$ giờ.\\ b) $\dfrac {40}{3}\,\mathrm {km/h}$.\\ Quãng đường $AB$ dài $30\,\mathrm {km}$. Tổng thời gian là $4{,}5$ giờ, tổng quãng đường là $60\,\mathrm {km}$. Vận tốc trung bình không bằng trung bình cộng của $12\,\mathrm {km/h}$ và $15\,\mathrm {km/h}$.  \hfill \qedEX 
\end{Solbook}
\begin{Solbook}{4}
 \par \noindent  $72$ học sinh.\\ Gọi số học sinh lớp $7$B lúc đầu là $5k$ ($k$ nguyên dương), số học sinh lớp $7$A là $4k$. Ta có $4k-8=\dfrac {1}{2}(5k+8)$. Suy ra $k=8$, nên $7$A có $32$ em và $7$B có $40$ em.  \hfill \qedEX 
\end{Solbook}
\begin{Solbook}{5}
 \par \noindent  a) Có. Vì $HD\parallel GE$ và $GE\parallel IF$ nên $HD\parallel IF$.\\ b) $x=65^{\circ }$, $y=115^{\circ }$.\\ $\widehat {HAB}$ và $\widehat {ABG}$ là hai góc trong cùng phía, nên $\widehat {ABG}=115^{\circ }$. Góc này kề bù với $\widehat {ABE}$, vậy $x=65^{\circ }$. Góc tương ứng của $x$ tại $C$ là $\widehat {BCF}=65^{\circ }$, suy ra $y=180^{\circ }-65^{\circ }=115^{\circ }$.  \hfill \qedEX 
\end{Solbook}
\begin{Solbook}{6}
 \par \noindent  a) $AB=BC=BD$ nên $AB=BD$. $\widehat {ABD}=\widehat {ABC}+\widehat {CBD}=120^{\circ }$. Tam giác $ABD$ cân tại $B$, hai góc đáy bằng $30^{\circ }$. Vậy $\widehat {BAD}=30^{\circ }$.\\ b) Tam giác $ABC$ đều, $H$ là trung điểm của $BC$ nên $AH$ là đường trung tuyến, đồng thời là đường cao. Vậy $AH\perp BC$.\\ c) $A$, $H$, $D$ thẳng hàng.\\ \textit {Hướng dẫn.} Trong tam giác đều $ABC$, $\widehat {BAH}=30^{\circ }$. Góc này bằng $\widehat {BAD}$ nên tia $AH$ trùng tia $AD$.  \hfill \qedEX 
\end{Solbook}
\begin{Solbook}{7}
 \par \noindent  $\sqrt {35}+\sqrt {99}<16$.\\ \textit {Hướng dẫn.} $\sqrt {35}<\sqrt {36}=6$ và $\sqrt {99}<\sqrt {100}=10$.  \hfill \qedEX 
\end{Solbook}
